\documentclass[11pt]{article}
\usepackage[margin=1in]{geometry}
\usepackage{amsmath,amssymb,amsthm,mathtools,microtype,needspace}
\usepackage[hidelinks]{hyperref}
\hypersetup{pdftitle={Fourfold forrelation: a self-contained proof},pdfauthor={Amir H. Safavi-Naeini}}
\newtheorem{theorem}{Theorem}
\newtheorem{lemma}[theorem]{Lemma}
\newtheorem{proposition}[theorem]{Proposition}
\newcommand{\E}{\mathbb E}
\newcommand{\Var}{\operatorname{Var}}
\newcommand{\TV}{\operatorname{TV}}
\newcommand{\sgn}{\operatorname{sgn}}
\newcommand{\diag}{\operatorname{diag}}
\newcommand{\wh}{\widehat}
\newcommand{\1}{\mathbf 1}
\DeclareMathOperator{\tr}{tr}
\title{Fourfold forrelation:\\a self-contained proof}
\author{Amir H. Safavi-Naeini}
\date{}
\begin{document}
\maketitle
\begin{abstract}
Four binary phase masks encode a scalar that a two-pass interferometer
can read directly. Three photons suffice to determine its promised sign
with error below one third, at total dose six. We prove that a single
parallel interrogation needs dose at least $(2/15)N^{1/8}$ for powers of
two $N\ge2^{30}$. The argument proceeds from the optical experiment to
parity functions, then constructs inputs whose small correlations defeat
low-dose measurements. All mathematical ingredients are proved here.
At $N=4096$ we also give an explicit finite criterion; the estimate needed
to turn that reduction into a dose-six lower bound remains open.
\end{abstract}

\section{The question and the answer}
Let $N=2^d$. Index modes by the binary vectors in $\mathbb F_2^d$ and put
\[
 H_{ij}=N^{-1/2}(-1)^{i\cdot j},\qquad
 u=N^{-1/2}(1,\ldots,1)^{\mathsf T}.
\]
The dot product in the exponent is modulo two. Thus $H$ is real,
symmetric and orthogonal: $H^2=I$. Each mask is a sign vector
$x^{(b)}\in\{\pm1\}^N$, represented by $D_b=\diag(x^{(b)})$.
The quantity of interest is
\begin{equation}\label{eq:forrelation}
 F(x)=u^{\mathsf T}D_1HD_2HD_3HD_4u.
\end{equation}
It lies in $[-1,1]$, since every matrix in the product is orthogonal.
For any vector $v$ we also write $D_v=\diag(v)$.
The input is promised to satisfy either $F\ge1/4$ or $F\le-1/4$.
The task is to decide which, with error at most $1/3$ on every promised input.

\paragraph{What is counted.}
A photon crossing one mask costs one unit of dose. A photon crossing
two masks costs two. A \emph{hard} dose bound holds on the entire support
of the signal state, including every branch of a superposition; it is
not an average. Arbitrary reference systems, called idlers, may bypass
the masks without cost.

In a \emph{parallel interrogation}, the whole input state is prepared
before any mask output is available. Each signal photon crosses the
direct-sum mask $D_1\oplus\cdots\oplus D_4$ once, and an arbitrary joint
measurement follows. Entanglement, repeated occupation of a mode, and
coherence between photon numbers are allowed. Jointly prepared time bins
are included; feeding an earlier output into the preparation of a later
probe is a different model.

\Needspace{12\baselineskip}
\begin{theorem}\label{thm:main}
For the promise problem above:
\begin{enumerate}
\item For every $N$, three photons making two mask crossings each have
hard dose six and worst-case error at most $81/256<1/3$.
\item For every power of two $N\ge2^{30}$, a parallel interrogation
with worst-case error at most $1/3$ has hard dose
\[
                         D\ge\frac{2}{15}N^{1/8}.
\]
\end{enumerate}
\end{theorem}

The second bound exceeds six for powers of two $N\ge2^{44}$.
It does not settle $N=4096$. Section~\ref{sec:finite} develops that
finite question separately and identifies the precise remaining estimate.
The main proof uses elementary linear algebra and probability; the
Gaussian tools used along the way are proved in Appendix~\ref{app:tools}.

\section{Why two passes are enough}
The product in \eqref{eq:forrelation} can be split into the inner product
of two unit vectors:
\[
 L_x=D_2HD_1u,\qquad R_x=HD_3HD_4u,\qquad
 \langle L_x,R_x\rangle=F(x).
\]
Prepare one photon in a uniform mode state and an equal superposition
of two paths, $(|0\rangle+|1\rangle)u/\sqrt2$.
Along path zero apply $D_1,H,D_2$; along path one apply $D_4,H,D_3,H$.
Both paths cross exactly two masks. The state becomes
\[
             \frac{|0\rangle L_x+|1\rangle R_x}{\sqrt2}.
\]
Recombine the paths and measure in the basis
$(|0\rangle\pm|1\rangle)/\sqrt2$, discarding the mode label.
The output port $X\in\{\pm1\}$ satisfies
\[
 \Pr(X=\pm1\mid x)=\frac{1\pm F(x)}2.
\]
On either promise, this flag is correct with probability at least $5/8$.
Repeat independently three times and take the majority. The probability
of at least two wrong flags is at most
\[
                 (3/8)^3+3(5/8)(3/8)^2=81/256<1/3.
\]
Every run uses dose six, and no outcome is discarded. This proves the
first part of Theorem~\ref{thm:main}.

\section{The lower-bound strategy}
A lower bound must limit \emph{every} allowed preparation and measurement.
We do this by finding two random input laws, $\mu_+$ and $\mu_-$, that
are hard to distinguish. They have two complementary properties:
\begin{itemize}
\item Each law places more than $9/10$ of its mass on its own promise.
Thus any correct solver must accept noticeably more often under $\mu_+$.
\item Products of a few input signs have almost the same expectation
under the two laws. A low-dose parallel experiment cannot collect enough
of these small differences to distinguish them.
\end{itemize}
The second point needs more than a bound on polynomial degree: there
may be many small correlations. The next section bounds their combined
effect on any measurement. Sections~\ref{sec:gaussian}--\ref{sec:closing}
then construct the laws, prove concentration, and obtain the contradiction
$1/5<1/12$.

\section{What a parallel measurement can learn}\label{sec:fourier}
\subsection{From photons to parity functions}
Write $m=4N$ for the total number of signal coordinates. An occupation
vector $\nu=(\nu_1,\ldots,\nu_m)$ with $n=\sum_j\nu_j\le D$ embeds
as the normalized symmetric sum of length-$D$ words containing $\nu_j$
copies of $j$ and $D-n$ copies of $0$. Distinct occupations give
orthogonal sums, so this is an isometry, including photon-number
superpositions. The extra symbol $0$ means an unused \emph{signal} slot;
all bypass modes belong to the unrestricted idler. Additional labels
on which the mask acts trivially, such as time-bin labels, can likewise
be retained in a slot's idler factor and then grouped with the other idlers.
The query acts as
\[
 O_x^{\otimes D},\qquad O_x|0\rangle=|0\rangle,\qquad
 O_x|j\rangle=x_j|j\rangle\quad(1\le j\le m).
\]
Allowing arbitrary slot states, beyond symmetric ones, only enlarges
the class of experiments against which we prove the bound.

For a word $a$, let $\partial a$ be the set of nonzero labels occurring
oddly. Repeated occurrences cancel in pairs because $x_j^2=1$.
The word's phase is therefore the parity function
$\chi_{\partial a}(x)$, where $\chi_S(x)=\prod_{j\in S}x_j$.
Purify the input and write
\[
 |\psi\rangle=\sum_a|a\rangle\otimes|\psi_a\rangle,\qquad
 \sum_a\|\psi_a\|^2=1.
\]
The idler vectors $\psi_a$ need not be orthogonal.
Let $f(x)$ be the probability of answering \emph{positive}. A measurement
effect $0\le E\le I$ gives this probability. Center it by putting
$T=E-I/2$, so $\|T\|\le1/2$. Then
\begin{equation}\label{eq:slot}
 f(x)-\tfrac12=\sum_{a,b}\chi_{\partial a\triangle\partial b}(x)
                  \langle\psi_a,T_{ab}\psi_b\rangle,\qquad
 T_{ab}=(\langle a|\otimes I)T(|b\rangle\otimes I).
\end{equation}
Here $\triangle$ is symmetric difference: labels common to both sets
cancel. The parity functions form an orthonormal basis for functions
on sign strings under the uniform measure. Consequently
\[
 f(x)=\sum_S\wh f(S)\chi_S(x),\qquad
 \wh f(S)=2^{-m}\sum_x f(x)\chi_S(x),
\]
and \eqref{eq:slot} shows $\wh f(S)=0$ when $|S|>2D$.

\subsection{A bound on an entire Fourier level}
We use one elementary matrix fact. If a scalar matrix factors as
$K_{ab}=\langle p_a,q_b\rangle$, with feature norms at most $P,Q$,
then for any block operator $T$,
\begin{equation}\label{eq:schur}
                   \|[K_{ab}T_{ab}]\|\le PQ\|T\|.
\end{equation}
Attach $q_b$ to input block $b$, apply $T$ without changing the attached
factor, and contract output block $a$ against $p_a$. The three operator
norms are at most $Q,\|T\|,P$. This also proves the claim for arbitrary
idler spaces. Here and below an unadorned norm is the Euclidean vector
norm or its induced operator norm.

\begin{lemma}[Fourier collection bound]\label{lem:fourier}
For $1\le r\le\min(2D,m)$ and arbitrary complex coefficients $c_S$,
\begin{equation}\label{eq:fourier}
 \left|\sum_{|S|=r}c_S\wh f(S)\right|
 \le\frac12\binom{2D}{r}m^{r/4}\max_{|S|=r}|c_S|.
\end{equation}
\end{lemma}
\begin{proof}
We partition the pairs $(a,b)$ in \eqref{eq:slot} with
$|\partial a\triangle\partial b|=r$. Each differing label occurs oddly
on exactly one side. Mark its \emph{last} occurrence on that side.
This assigns a unique record of $r$ marked positions to the pair.

For example, $a=(1,1,2,3)$ and $b=(2,4,0,0)$ have parity sets
$\{2,3\}$ and $\{2,4\}$. Mark the $3$ in $a$ and the $4$ in $b$.
Removing these marks leaves the same parity set $\{2\}$ on both sides.
The repeated $1$ and the unused slots contribute no parity.

Fix a record with $k$ marks in $a$ and $r-k$ in $b$.
A word is valid on its own side if the marked labels are nonzero,
distinct, occur oddly, and are marked at their last occurrences.
Write $A,B$ for the marked sets and
$R_a=\partial a\setminus A$, $R_b=\partial b\setminus B$ for the
residual sets. The pair belongs to the record exactly when both words
are valid, $A\cap B=\varnothing$, and $R_a=R_b$.
Indeed, then $\partial a=R_a\mathbin{\dot\cup}A$ and
$\partial b=R_a\mathbin{\dot\cup}B$, so precisely the marked labels differ.

Index an auxiliary matrix $C$ by $k$-element sets $A$ and $(r-k)$-element
sets $B$, putting $C_{AB}=c_{A\cup B}$ if they are disjoint and zero
otherwise. Factor it using standard basis vectors on its shorter side.
For example, if there are fewer rows, use a basis vector for each row
and the corresponding column of $C$ for each column. The product of
the largest feature norms is at most
\[
 \sqrt{\min\left\{\binom{m}{k},\binom{m}{r-k}\right\}}\max|c_S|
 \le m^{\min(k,r-k)/2}\max|c_S|.
\]
Attach to each feature the basis vector labelled by its residual set,
and set the feature to zero for an invalid word. Their inner product
now enforces residual equality and validity, giving the \emph{exact}
record matrix $K$, with the same norm bound.
By \eqref{eq:schur}, the absolute value of its contribution
$\langle\psi,[K_{ab}T_{ab}]\psi\rangle$ is at most half that bound.
There are $\binom{D}{k}\binom{D}{r-k}$ records with this split. Hence
\[
 \left|\sum_{|S|=r}c_S\wh f(S)\right|
 \le\frac12\max|c_S|
       \sum_{k=0}^r\binom{D}{k}\binom{D}{r-k}m^{\min(k,r-k)/2}.
\]
Use $\min(k,r-k)\le r/2$ and count all ways to choose $r$ positions
among $2D$ slots to obtain \eqref{eq:fourier}.
\end{proof}

The lemma converts small input moments into a bound on any experiment.
If two laws have moment differences
$\delta_r=\max_{|S|=r}|(\E_+-\E_-)\chi_S|$, then
\begin{equation}\label{eq:collection}
 |\E_+f-\E_-f|
 \le\frac12\sum_{r=1}^{\min(2D,m)}\binom{2D}{r}m^{r/4}\delta_r.
\end{equation}
The constant Fourier coefficient cancels because both laws have mass one.

\section{Inputs with small correlations}\label{sec:gaussian}
We want four sign blocks with appreciable forrelation but very small
low-order moments. Three independent correlated links supply the
forrelation. Random block signs make every distinguishing moment use
all four blocks.

Fix $\rho=200/201$. Independently for $s=1,2,3$, draw jointly Gaussian
standard vectors $Z_s,W_s$ with
$\E Z_sW_s^{\mathsf T}=\rho H$, and set
$a_s=\sgn Z_s$, $b_s=\sgn W_s$ componentwise.
This covariance is valid because $\|\rho H\|<1$.
Join the links by componentwise multiplication, denoted $\odot$:
\[
          y=(a_1,\ b_1\odot a_2,\ b_2\odot a_3,\ b_3).
\]
To define $\mu_\sigma$ for $\sigma\in\{\pm1\}$, choose independent
uniform signs $\zeta_1,\zeta_2,\zeta_3$, set
$\zeta_4=\sigma\zeta_1\zeta_2\zeta_3$, and put $x^{(b)}=\zeta_b y^{(b)}$.
Since $F$ is linear in each block, $F(x)=\sigma F(y)$.

\subsection{A single link}
For coordinate sets $A,B$ write
$\mathcal L(A,B)=\E\prod_{i\in A}a_i\prod_{j\in B}b_j$.
Let $a=|A|$, $b=|B|$, and suppose $a+b\le r\le\sqrt N$.
Then
\begin{equation}\label{eq:link}
 |\mathcal L(A,B)|\le
       \left(\frac{r}{2\sqrt N}\right)^{\max(a,b)}.
\end{equation}
If just one set is empty, its moment is zero because the other marginal
has independent symmetric signs; if both are empty, the moment is one.
Otherwise the cross-covariance $R=\rho H_{A,B}$ has
\[
 \|R\|\le\left(\sum_{i,j}|R_{ij}|^2\right)^{1/2}=\rho\sqrt{ab/N}
                         \le r/(2\sqrt N)\le1/2.
\]
A product of signs in $a$ independent Gaussian coordinates has no
Hermite terms of degree below $a$: it is odd in each coordinate.
Here Hermite expansion means orthogonal polynomial expansion for Gaussian inputs.
The same holds with $b$ on the other side. The Gaussian contraction
proved in Appendix~\ref{app:tools} pairs only equal degrees $p$ and
contracts degree $p$ by at most $\|R\|^p$. Both sign products have
$L^2$ norm one. Their common degrees start at $\max(a,b)$, proving
\eqref{eq:link}.

\subsection{Three links and four odd blocks}
For $S=(S_1,\ldots,S_4)$ put $q_b=|S_b|$ and $r=\sum_bq_b$.
Averaging the block signs gives
\[
 \E_\zeta\prod_{b=1}^4\zeta_b^{q_b}
 =\sigma^{q_4}\prod_{j=1}^3\1_{\{q_j\equiv q_4\pmod2\}}.
\]
The difference between the two laws therefore vanishes unless all
four $q_b$ are odd. In that case independence of the links gives
\[
 (\E_{\mu_+}-\E_{\mu_-})\chi_S
       =2\mathcal L(S_1,S_2)\mathcal L(S_2,S_3)\mathcal L(S_3,S_4).
\]
The three exponents in \eqref{eq:link} have sum at least $3r/4$:
\[
 \begin{aligned}
 \max(q_1,q_2)&\ge(3q_1+q_2)/4,\\
 \max(q_2,q_3)&\ge(q_2+q_3)/2,\\
 \max(q_3,q_4)&\ge(q_3+3q_4)/4.
 \end{aligned}
\]
Because the base in \eqref{eq:link} is at most one, we conclude
\begin{equation}\label{eq:moments}
 \delta_r\le2^{1-3r/4}r^{3r/4}N^{-3r/8}\quad(r\le\sqrt N),
 \qquad \delta_r=0\quad\text{unless $r\ge4$ is even}.
\end{equation}
The scale of the final result is already visible: at degree four the
moment is $O(N^{-3/2})$, while the Fourier bound costs $O(D^4N)$.
Their product is $O(D^4/\sqrt N)$, which becomes constant at $D=N^{1/8}$.

\section{Why these inputs satisfy the promise}\label{sec:concentration}
Small moments alone do not make hard inputs: the laws must also place
most of their mass on the required promises. We prove this by computing
the mean and bounding the variance.

\subsection{The mean is beyond the threshold}
The Gaussian sign identity proved in Appendix~\ref{app:tools} gives
\[
 \E a_i b_j=\beta_N H_{ij},\qquad
 \beta_N=\frac{2\sqrt N}{\pi}\arcsin(\rho/\sqrt N).
\]
Expanding \eqref{eq:forrelation} and using the three independent links,
\[
 \E_{\mu_\sigma}F
 =\frac{\sigma\beta_N^3}{N}\sum_{i,j,k,l}H_{ij}^2H_{jk}^2H_{kl}^2
 =\sigma\beta_N^3.
\]
The sum has $N^4$ terms, each equal to $N^{-3}$.
Since $\arcsin t\ge t$ for $t\ge0$ and $\pi<22/7$,
\begin{equation}\label{eq:margin}
 \beta_N^3-\frac14>
       \left(\frac{1400}{2211}\right)^3-\frac14>\frac1{260}.
\end{equation}
The final comparison is the integer inequality
$130\cdot1400^3-33\cdot2211^3=38\,842\,277>0$.

\subsection{Separate independent noise from smooth variation}
\begin{lemma}\label{lem:variance}
Under either law, $\Var F<1540/N$.
\end{lemma}
\begin{proof}
Realize each link using independent standard Gaussian vectors
$G,\xi,\eta$ as
\[
 a=\sgn(G+\xi/\sqrt{200}),\qquad
 b=\sgn(HG+\eta/\sqrt{200}).
\]
After standardizing the marginal variances, the cross-covariance is
$\rho H$, as required. Condition on $\mathcal G=(G_1,G_2,G_3)$.
All primitive signs are then independent, and each is used in only
one coordinate of $y$. Hence the coordinates of $y$ are conditionally
independent too. Their primitive means are $\phi(G)$ and $\phi(HG)$,
where
\[
 \phi(t)=2\Phi(\sqrt{200}\,t)-1=\operatorname{erf}(10t),\qquad
 L_\phi=\|\phi'\|_\infty=20/\sqrt\pi,\qquad L_\phi^2<128.
\]
Here $\Phi$ is the standard normal distribution function; the last
bound uses $\pi>25/8$. Since $F(x)=\sigma F(y)$, the gauges do not
affect its variance. We use the two terms of total variance:
\begin{equation}\label{eq:total-variance}
 \Var F=\E\Var(F\mid\mathcal G)+\Var\overline F,\qquad
 \overline F=\E(F\mid\mathcal G).
\end{equation}

\emph{First term: independent sign noise.}
For independent signs, variance tensorization bounds the variance of
a multilinear function by the sum of its expected squared coefficients
in each sign; Appendix~\ref{app:tools} gives the short proof.
Each endpoint block of $F$ contributes exactly $1/N$ to this sum:
its other chain has norm one, and the endpoint vector is flat.
For block two the coefficient of coordinate $j$ is $\ell_jv_j$, where
\[
 \ell=HD_{a_1}u,\qquad v=HD_{b_2\odot a_3}HD_{b_3}u.
\]
These vectors depend on disjoint links and are independent.
The coordinates of $a_1$ are independent uniform signs, so
$\E\ell_j^2=1/N$; also $\|v\|=1$. Averaging the sum gives $1/N$.
Block three has the same bound by reversing the argument, using the
independent uniform signs in $b_3$. Consequently
\begin{equation}\label{eq:noise-var}
                         \E\Var(F\mid\mathcal G)\le4/N.
\end{equation}

\emph{Second term: a smooth Gaussian function.}
Set
\[
 M(g)=D_{\phi(g)}H D_{\phi(Hg)}.
\]
This is a contraction because $|\phi|\le1$, and the conditional mean is
$\overline F=\sigma u^{\mathsf T}M(G_1)M(G_2)M(G_3)u$.
We first establish a useful symmetry of partial chains.
For a binary shift $s$ define
$(P_sv)_i=v_{i+s}$ and $(C_sv)_i=(-1)^{s\cdot i}v_i$.
The identities $HP_s=C_sH$, $HC_s=P_sH$, and oddness of $\phi$ give
\begin{equation}\label{eq:translation}
                  M(P_sg)=P_sM(g),\qquad M(C_sg)=M(g)P_s.
\end{equation}
Both $P_sG$ and $C_sG$ have the same law as $G$.
Thus any vector obtained by applying an independent chain of $M$'s
or transposed $M$'s to $u$ has translation-invariant distribution:
change the Gaussian in its outermost factor using \eqref{eq:translation}.
Its norm is at most one. All coordinate second moments are equal, so
\begin{equation}\label{eq:flatness}
                    \E|v_i|^2=\frac1N\E\|v\|^2\le\frac1N.
\end{equation}
The empty chain $u$ obeys the same bound.

At one of the three links, write the conditional mean, up to $\sigma$,
as $p^{\mathsf T}M(G)v$. The left and right partial chains $p,v$ and
the current $G$ are mutually independent; both chains satisfy
\eqref{eq:flatness}. Differentiating this smooth expression and using
$\|a+b\|^2\le2\|a\|^2+2\|b\|^2$ gives
\[
 \|\nabla_G\overline F\|^2\le2L_\phi^2
 \left(\|p\odot HD_{\phi(HG)}v\|^2+
       \|v\odot HD_{\phi(G)}p\|^2\right).
\]
For the first norm, condition on $G,v$ and average over $p$;
\eqref{eq:flatness} gives
\[
 \E_p\|p\odot HD_{\phi(HG)}v\|^2
       \le\frac1N\|HD_{\phi(HG)}v\|^2\le\frac1N.
\]
The other norm is symmetric.
The Gaussian Poincare inequality, proved in Appendix~\ref{app:tools},
now gives
\[
 \Var\overline F\le\sum_{s=1}^3\E\|\nabla_{G_s}\overline F\|^2
                         \le12L_\phi^2/N.
\]
Combining this with \eqref{eq:noise-var} in \eqref{eq:total-variance}
yields $\Var F\le(4+12L_\phi^2)/N<1540/N$.
\end{proof}

\section{Closing the lower bound}\label{sec:closing}
The mean margin \eqref{eq:margin}, Lemma~\ref{lem:variance}, and
Chebyshev's inequality imply, for $N\ge2^{30}$,
\begin{equation}\label{eq:failure}
 \varepsilon_\sigma:=\mu_\sigma\{\sigma F<1/4\}
 \le\frac{1540\cdot260^2}{N}
 =\frac{104\,104\,000}{N}<\frac1{10}.
\end{equation}
The last step follows from $1\,041\,040\,000<2^{30}$.
A correct acceptance probability $f\in[0,1]$ obeys $f\ge2/3$
on the positive promise and $f\le1/3$ on the negative promise.
Even allowing the worst behavior off the promises,
\[
 \E_{\mu_+}f\ge\tfrac23(1-\varepsilon_+),\qquad
 \E_{\mu_-}f\le\tfrac13(1-\varepsilon_-)+\varepsilon_-.
\]
Thus correctness requires
\begin{equation}\label{eq:required-gap}
 \E_{\mu_+}f-\E_{\mu_-}f
 \ge\frac13-\frac23(\varepsilon_++\varepsilon_-)>\frac15.
\end{equation}

Suppose instead $D<(2/15)N^{1/8}$. Then $2D<\sqrt N$, so all
relevant moments satisfy \eqref{eq:moments}.
Substitute that bound and $m=4N$ into \eqref{eq:collection}, using
$\binom{2D}{r}\le(2eD/r)^r$:
\[
 |\E_{\mu_+}f-\E_{\mu_-}f|
 \le\sum_{\substack{4\le r\le2D\\r\text{ even}}}
        \left(2^{3/4}eD\,r^{-1/4}N^{-1/8}\right)^r.
\]
For $r\ge4$ each base is at most
$\theta=2^{1/4}eDN^{-1/8}<1/2$, since $e<3$ and $2^{1/4}<5/4$.
The even-level sum is therefore at most
\[
             \sum_{j=2}^{\infty}\theta^{2j}
             =\frac{\theta^4}{1-\theta^2}<\frac1{12}.
\]
This contradicts \eqref{eq:required-gap} and proves the second part
of Theorem~\ref{thm:main}. The argument has not conditioned the laws:
accounting directly for their small off-promise mass is sufficient.

\section{The finite question at \texorpdfstring{$N=4096$}{N=4096}}\label{sec:finite}
Can a parallel probe of dose six solve the same task at $N=4096$?
The asymptotic estimate is too weak to answer this. We now give a finite
construction and an explicit reduction of the proposed obstruction to a
matrix estimate. The reduction is proved; the matrix estimate is open.
All its coefficients are defined below, so the question has no implicit
numerical or documentary inputs.

\subsection{An exact plant}
Let $q=64$, $N=q^2$, and identify a coordinate with a pair $(a,b)$ of
binary vectors of length six. Write $S_q=\sqrt q\,H_q$, so its entries
are signs and $H_N=H_q\otimes H_q$.
For a uniform permutation $\pi$ of these $q$ labels and independent
uniform signs $\sigma_b$, define arrays
\begin{equation}\label{eq:finite-link}
 L_{a,b}=S_q(a,\pi(b))\sigma_b,\qquad
 R_{\pi(b),a}=\sigma_bS_q(b,a).
\end{equation}
They obey $H_NL=R$ and $H_NR=L$. Indeed, with $b_c=\pi^{-1}(c)$,
\[
 (H_NL)_{c,d}
 =\frac1q\sum_{a,b}S_q(c,a)S_q(d,b)S_q(a,\pi(b))\sigma_b
 =S_q(d,b_c)\sigma_{b_c}=R_{c,d};
\]
the other identity follows from $H_N^2=I$.
Take three independent links and form
\begin{equation}\label{eq:plant}
                z=(L_1,R_1\odot L_2,R_2\odot L_3,R_3).
\end{equation}
Starting on the right of \eqref{eq:forrelation}, the transform identities
successively cancel $R_3,L_3,R_2,L_2,R_1,L_1$ and leave $u$.
Thus $F(z)=1$ exactly; negating the first block gives $F=-1$.

Independently multiply every coordinate by a sign of mean $\beta=19/25$,
equivalently flipping it with probability $3/25$.
Call the resulting positive law $\widetilde P_+$ and its first-block
reflection $\widetilde P_-$. Their means are $\pm\beta^4$.
Condition each law onto its respective promise to obtain $P_+,P_-$.
For probability laws $P,Q$, we use
$\TV(P,Q)=\sup_A|P(A)-Q(A)|$, half their $L^1$ distance.

\begin{lemma}[Cost of conditioning]\label{lem:conditioning}
Let
\begin{equation}\label{eq:kappa}
 V_\beta=\frac{(1+\beta^2)(1+\beta^4)}N,\qquad
 g=\beta^4-\frac14,\qquad
 \kappa=2\exp\left(-\frac{g^2}{2V_\beta}\right).
\end{equation}
Conditioning changes the output total variation of any experiment
by at most $\kappa$, and $\kappa<1/400$.
\end{lemma}
\begin{proof}
Fix an exact plant $z$ and reveal the four noisy blocks from left to
right. The conditional expectations of $F$ form a martingale.
At block $b$, averaging future noise contributes $\beta^{4-b}$.
The remaining deterministic right chain is a sign vector divided by
$\sqrt N$, by the plant identities, and the left chain has norm one.
Thus the martingale increment is a sum of independent centered signs
with sum of squared coefficients $\beta^{2(4-b)}/N$.

A centered variable in an interval of length two satisfies
$\E e^{tY}\le e^{t^2/2}$: for its log moment-generating function,
the second derivative is a tilted variance at most one, and its value
and first derivative vanish at zero. Apply this conditionally to each
sign and then each block. The four increments give
$\E e^{t(F-\beta^4)}\le e^{t^2V_\beta/2}$.
Markov's inequality, optimized at $t=-g/V_\beta$, bounds the promise
failure by $\exp(-g^2/(2V_\beta))$. Reflection gives the same bound
on the negative side.

Conditioning a law on an event of probability $1-\epsilon$ changes
it in total variation by exactly $\epsilon$ (split the law into its
on-event and off-event parts). An experiment cannot increase total
variation. The triangle inequality gives the sum of the two failures,
which is \eqref{eq:kappa}.
For an elementary numerical bound, direct rational substitution gives
\[
 \frac{g^2}{2V_\beta}
   =\frac{546\,296\,776\,992}{80\,258\,243\,125}>\frac{34}{5}.
\]
Also $e>27/10$ from its positive power series, and
$e^{4/5}>1+4/5+(4/5)^2/2=53/25$. Hence
$e^{34/5}>(27/10)^6(53/25)>800$, proving $\kappa<1/400$.
\end{proof}

\subsection{The exact matrix seen by a parallel probe}
For a Fock occupation $\nu$, define its parity support in each block by
$S_b(\nu)=\{i:\nu_{b,i}\text{ is odd}\}$.
The mask acts only through this support:
$U_x|\nu\rangle=\chi_{S(\nu)}(x)|\nu\rangle$.
For hard dose six, $|S|:=\sum_b|S_b|\le6$.
Purify any input and group all terms with support $S$ into a vector
$\psi_S$, including their idlers. These vectors are mutually orthogonal,
because they contain disjoint collections of occupation-basis vectors.
Set $p_S=\|\psi_S\|^2$.

For the unconditioned laws define
\begin{equation}\label{eq:finite-K}
 K_{ST}=\tfrac12(\E_{\widetilde P_+}-\E_{\widetilde P_-})
                                \chi_{S\triangle T}.
\end{equation}
The map $|S\rangle\mapsto\psi_S/\sqrt{p_S}$ is an isometry for nonzero
weights. Therefore half the difference of the two averaged purified
output states has trace norm
\begin{equation}\label{eq:finite-trace}
                   \|\diag(\sqrt p)K\diag(\sqrt p)\|_1.
\end{equation}
The trace norm is the sum of singular values. For a trace-zero Hermitian
matrix $X$, its positive and negative parts have equal trace, giving
$\sup_{0\le E\le I}|\tr(EX)|=\|X\|_1/2$.
Thus \eqref{eq:finite-trace} bounds every measurement's output total
variation, even after a purifier is discarded. Mixtures are covered
by purification as well.

Globally negating both halves $(L_j,R_j)$ of one link preserves its law.
Consequently a plant moment vanishes unless adjacent block degrees
have the same parity. First-block reflection in \eqref{eq:finite-K}
selects the odd case: only differences with an odd number of labels in
\emph{every} block survive. Independent noise multiplies a degree-$r$
moment by $\beta^r$.

\subsection{From mode labels to 210 count states}
Here is a completely specified finite criterion. For nonnegative
integer four-vectors $s,t$ with $|s|,|t|\le6$, let $A_b,B_b$ range over sets of sizes
$s_b,t_b$. Define the complete kernel and its coefficient by
\begin{align}
 \mathcal O_{s,t}(A,B)
   &=\1_{\{A_b\cap B_b=\varnothing\ \text{for all }b\}}
                      \E_z\chi_{A\cup B}(z),\label{eq:complete}\\
 c_{s,t}
   &=\sup_{\alpha,\omega}
       \|\diag(\sqrt\alpha)\mathcal O_{s,t}
                           \diag(\sqrt\omega)\|_1.\label{eq:cstar}
\end{align}
The expectation is the finite uniform average over the three independent
permutations and sign families in \eqref{eq:finite-link}; thus it is
explicitly defined. The supremum is over all probability laws on the
row and column sets, including correlated laws between blocks.
These are finite matrices and compact probability simplexes, so the
suprema exist. Transposition gives $c_{s,t}=c_{t,s}$.
The disjointness indicator is part of the matrix, not a removable
qualification on selected entries.

Write $|n|=\sum_bn_b$ and
$\binom nr=\prod_{b=1}^4\binom{n_b}{r_b}$. Index a symmetric
nonnegative matrix $B$ by
\[
             \mathcal N=\{n\in\mathbb Z_{\ge0}^4:|n|\le6\},
             \qquad |\mathcal N|={10\choose4}=210.
\]
For $n,m\in\mathcal N$, set
\begin{equation}\label{eq:B}
 B_{nm}=\1_{\{n_b+m_b\ \mathrm{odd}\ \forall b\}}
 \sum_{0\le r\le\min(n,m)}
 \beta^{|n|+|m|-2|r|}c_{n-r,m-r}
                         \sqrt{\binom nr\binom mr}.
\end{equation}
Inequalities between four-vectors are componentwise.

\begin{proposition}[Finite spectral reduction]\label{prop:finite}
For every parallel probe of hard dose at most six, including arbitrary
idlers, repeated occupations and photon-number coherence, the output
laws under $P_+,P_-$ satisfy
\begin{equation}\label{eq:finite-goal}
                         \TV(T_+,T_-)\le\lambda_{\max}(B)+\kappa.
\end{equation}
In particular, $\lambda_{\max}(B)+\kappa<1/3$ would rule out a
worst-case error of $1/3$ at dose six.
\end{proposition}
\begin{proof}
First fix the block counts $n_b=|S_b|$, $m_b=|T_b|$ and the common
support $R=S\cap T$, with counts $r_b=|R_b|$.
Write $S=R\mathbin{\dot\cup}A$, $T=R\mathbin{\dot\cup}B'$.
The exclusive counts are $s=n-r$, $t=m-r$.
The corresponding block of $K$ is the kernel \eqref{eq:complete},
restricted to sets avoiding $R$, multiplied by $\beta^{|s|+|t|}$.
It vanishes unless $n_b+m_b$ is odd for every block.

To bound this block, temporarily take normalized laws $p_S,w_T$ on
the two count sectors. For a fixed $R$ put
$P_R=\sum_Ap_{R\cup A}$ and $W_R=\sum_{B'}w_{R\cup B'}$.
The definition \eqref{eq:cstar}, applied to the restricted normalized
laws, bounds the weighted trace norm by
$c_{s,t}\sqrt{P_RW_R}$ before the noise factor.
Each nonzero entry has the unique residual $R=S\cap T$, so summing
over all residuals reconstructs the desired block exactly.
Triangle inequality and Cauchy--Schwarz give
\[
 \sum_Rc_{s,t}\sqrt{P_RW_R}
 \le c_{s,t}\sqrt{\binom nr\binom mr}.
\]
Indeed each $S$ contains exactly $\binom nr$ eligible residual sets,
and each $T$ exactly $\binom mr$.

Now let $v_n=\sqrt{\sum_{S:\,|S_b|=n_b}p_S}$ for the actual input.
Restoring sector masses supplies a factor $v_nv_m$. Sum over common
counts $r$ and ordered sectors $n,m$, using triangle inequality again.
Equations~\eqref{eq:finite-trace} and \eqref{eq:B} give
\[
 \TV(\widetilde T_+,\widetilde T_-)
       \le v^{\mathsf T}Bv\le\lambda_{\max}(B),
       \qquad \|v\|_2=1.
\]
Lemma~\ref{lem:conditioning} adds at most $\kappa$.
Finally, equal-prior optimal decision error is $(1-\TV)/2$: choose
the more likely law for each output and sum the pointwise minima.
A worst-case error at most $1/3$ would give equal-prior error at most
$1/3$, requiring $\TV\ge1/3$.
\end{proof}

\subsection{A worked example: dose two}
The criterion does give a concrete finite conclusion at smaller dose.
With hard dose two, any surviving difference uses all four blocks, so
the two supports must be disjoint, of size two each, and occupy
complementary pairs of blocks. For coordinates $i=(a,b)$, $j=(c,d)$,
sign averaging in \eqref{eq:finite-link} leaves only $\pi(b)=c$; hence
\[
 \E L_iR_j=\frac1q S_q(a,c)S_q(b,d)=H_N(i,j).
\]
The three independent links therefore make each entry of a surviving
kernel a product of three Walsh entries, of magnitude $N^{-3/2}$.
Each row has $N^2$ columns, so its squared norm is $1/N$.
For arbitrary probability laws $\alpha,\omega$, factor the weighted
matrix as $(\diag(\sqrt\alpha)\mathcal O)\diag(\sqrt\omega)$.
The inequality $\|XY\|_1\le\|X\|_{\mathrm F}\|Y\|_{\mathrm F}$, where
the Frobenius norm is the square root of the sum of squared entries,
gives $c_{s,t}\le1/\sqrt N$. This inequality follows by duality:
for $\|U\|\le1$, apply Cauchy--Schwarz to
$\tr(UXY)$ and use $\|UX\|_{\mathrm F}\le\|X\|_{\mathrm F}$.

Restricting the count states to $|n|\le2$ leaves three complementary
two-state blocks and zero rows. Each edge is at most $\beta^4/\sqrt N$.
Thus, for every parallel probe of dose at most two at $N=4096$,
\[
 \TV(T_+,T_-)\le\frac{\beta^4}{64}+\kappa
       <\frac{192821}{25000000}<\frac1{100}.
\]
Its equal-prior error exceeds $99/200$. This is a proved finite
obstruction; the dose-six problem requires higher-degree kernel bounds.

\subsection{What remains at dose six}
The entries of $B$ are defined entirely by finite averages and finite
norm optimizations. No bound $\lambda_{\max}(B)+\kappa<1/3$ is proved
here. Establishing it requires estimates for the \emph{complete} kernels
\eqref{eq:complete}; accurate numerical arithmetic on proposed upper
bounds cannot supply a missing analytic inequality.
This matrix estimate is sufficient, not necessary: failure to prove it,
or even its failure, would not exclude a sharper lower-bound argument.

One sometimes imposes an extra restriction that makes this calculation
smaller. Define the parity-size operator
\[
                 Q|\nu\rangle=|S(\nu)|\,|\nu\rangle.
\]
Extend $Q$ by the identity on idlers. For an input density operator
$\varrho$ with $[\varrho,Q]=0$,
decompose into fixed-$Q$ sectors; only pairs with
$|n|=|m|$ are needed. In \eqref{eq:B} this means replacing $B$ by
$B^{\rm bal}_{nm}=\1_{\{|n|=|m|\}}B_{nm}$. The same proof applies
sector by sector and then by convexity. Even the corresponding smaller
matrix estimate remains open.

Fixed physical photon number does not justify this restriction.
For example, three photons in coordinate zero of block one have parity
size one. One photon in coordinate zero of each of blocks two, three
and four has parity size three. Their superposition has fixed physical
number three but is not diagonal in $Q$. Its connecting moment is
nonzero: each link has $\E L_{0,0}R_{0,0}=1/q$, because sign averaging
leaves precisely $\pi(0)=0$. Thus \eqref{eq:finite-K} for these two
supports equals $\beta^4/q^3>0$. Deleting it would change the experiment.
Appendix~\ref{app:obstruction} gives another explicit obstruction to
a shortcut in estimating the finite coefficients.

\section{The remaining adaptive question}
The proved lower bound concerns a fixed parallel input and measurement.
If later probes depend on earlier outcomes, the common fixed operator
used in \eqref{eq:slot} is no longer available. For fresh batches with
classical feedback and total hard dose $D$, the acceptance probability
still has degree at most $2D$: along a fixed history each unnormalized
amplitude contains at most $D$ mask-sign factors. Its probability has
degree at most $2D$, and summing histories preserves this bound.
Degree alone is insufficient.
A sufficient additional estimate would be
\[
 \left|\sum_{|S|=r}c_S\wh f(S)\right|
 \le \binom{2D}{r}(4N)^{r/4}\max|c_S|.
\]
This is twice the parallel bound. It would give a gap below $1/6$ in
Section~\ref{sec:closing}, still contradicting the required gap above
$1/5$. We have not established this estimate for adaptive preparations.
Neither a mean-dose bound nor optimality of the dose-six protocol is
claimed.

\appendix
\section{The probability tools used in the proof}\label{app:tools}
The following short derivations make the Gaussian and concentration
steps independent of external references.

\subsection{The Gaussian sign identity}
If standard Gaussians $X,Y$ have correlation $c$, write
$X=G_1$, $Y=cG_1+\sqrt{1-c^2}\,G_2$ with independent standard
$G_1,G_2$. The angle in this plane is uniform. The two signs disagree
on wedges of total angle $2\arccos c$, so
\[
 \Pr(\sgn X\ne\sgn Y)=\frac{\arccos c}{\pi},\qquad
 \E\sgn X\sgn Y=\frac2\pi\arcsin c.
\]
The endpoint correlations follow by continuity.

\subsection{Hermite degrees and Gaussian contraction}
Define the normalized Hermite polynomials by
\[
 e^{tx-t^2/2}=\sum_{n\ge0}h_n(x)\frac{t^n}{\sqrt{n!}}.
\]
Taking Gaussian expectations of two generating functions shows
$\E h_n(G)h_m(G)=\1_{\{n=m\}}$.
The products $h_\alpha(x)=\prod_i h_{\alpha_i}(x_i)$ form an
orthonormal basis in Gaussian $L^2$. One way to see completeness is
to take a function orthogonal to all polynomials. Its Gaussian-weighted
Laplace transform has all derivatives zero at zero and is entire by
Cauchy--Schwarz and Gaussian exponential integrability. It therefore
vanishes also at imaginary arguments; uniqueness of the Fourier
transform makes the original function zero.

The total degree of $h_\alpha$ is $|\alpha|=\sum_i\alpha_i$.
The degree-$p$ space is invariant under orthogonal rotations, since
it is the orthogonal complement of degrees at most $p-1$ inside the
polynomials of degree at most $p$. If standard Gaussian vectors $X,Y$
have cross-covariance $R$, their generating functions satisfy
\[
 \E e^{t\cdot X-\|t\|^2/2}e^{s\cdot Y-\|s\|^2/2}=e^{t^{\mathsf T}Rs}.
\]
Rotate $X,Y$ by a singular-value decomposition of $R$ (padding to the
same dimension with independent coordinates if necessary).
If its singular values are $s_j$, expansion in these rotated coordinates gives
\[
 \E h_\alpha(X)h_\gamma(Y)
      =\1_{\{\alpha=\gamma\}}\prod_j s_j^{\alpha_j}.
\]
Thus different total degrees do not pair, and degree $p$ contracts by
at most $\|R\|^p$. If $f,g$ have no degrees below $a,b$, then, writing
$f_p,g_p$ for their degree-$p$ parts and using $\|R\|\le1$,
\[
 |\E f(X)g(Y)|
 \le\sum_{p\ge\max(a,b)}\|R\|^p\|f_p\|_2\|g_p\|_2
 \le\|R\|^{\max(a,b)}\|f\|_2\|g\|_2.
\]
A product of signs in $a$ distinct coordinates is odd in every one
of them. Each participating Hermite index is therefore odd, giving
minimum degree $a$ and $L^2$ norm one, as used in \eqref{eq:link}.

\subsection{Gaussian Poincare inequality}
Let $f=\sum_\alpha c_\alpha h_\alpha$ be smooth with square-integrable
gradient. Gaussian integration by parts gives
\[
 \langle\partial_i f,h_\gamma\rangle
     =\sqrt{\gamma_i+1}\,c_{\gamma+e_i}.
\]
Parseval's identity therefore yields
\[
 \Var f=\sum_{\alpha\ne0}|c_\alpha|^2
       \le\sum_\alpha|\alpha|\,|c_\alpha|^2
       =\E\|\nabla f\|^2.
\]
The smooth functions used in the proof have bounded derivatives, so
the integrations are justified by Gaussian integrability.

\subsection{Variance tensorization for independent variables}
For independent $Y_1,\ldots,Y_k$, let $E_i$ average coordinate $i$
alone and let $F_i=\E[F\mid Y_1,\ldots,Y_i]$.
The martingale differences are orthogonal, and independence gives
\[
 F_i-F_{i-1}=\E[F-E_iF\mid Y_1,\ldots,Y_i].
\]
Conditional Jensen and summation imply
\[
 \Var F\le\sum_i\E(F-E_iF)^2
           =\sum_i\E\Var_{Y_i}(F\mid Y_{-i}).
\]
If $Y_i$ is a sign and $F=A+Y_iB$ with $A,B$ independent of that
coordinate, the last conditional variance is
$(1-(\E Y_i)^2)B^2\le B^2$. This proves the coefficient bound used
in \eqref{eq:noise-var}.

\section{Why a finite coefficient shortcut fails}\label{app:obstruction}
This example explains why the finite matrix estimate is a substantive
remaining task. For the link \eqref{eq:finite-link}, define
\[
 M(A,B)=\E\prod_{v\in A}L_v\prod_{w\in B}R_w.
\]
For any power of two $q\ge8$, use the three-cell sets
\[
 V=\{(0,0),(1,0),(2,0)\},\qquad
 W=\{(0,0),(0,1),(0,2)\}.
\]
Here $0,1,2$ denote their distinct binary labels.
For $M(V,W)$, the sign in column zero must match the sign in
column $\pi^{-1}(0)$, so averaging leaves exactly $\pi(0)=0$.
Every remaining Walsh phase is $+1$, giving $M(V,W)=1/q$.
For $M(W,V)$, sign averaging requires the permutation to map
$\{0,1,2\}$ onto itself; again all surviving phases are $+1$.
Thus
\[
             M(V,W)=\frac1q,\qquad
             M(W,V)=\frac1{\binom q3}.
\]
The exact four-block plant moment at $(V,W,V,W)$ is consequently
$1/(q^2\binom q3)$. In particular, a proposed universal maximum
of the form
\[
 \frac{b^2}{\binom q3},\qquad
 b=\frac{q+2}{q(q-1)(q-2)},
\]
is false: $b<1/q$ for $q\ge8$.
At $q=8$ the actual entry is $1/3584$, exceeding the proposed
$25/1580544$. At $q=64$ the same comparison applies.
This refutes that entrywise shortcut; it does not refute a possible
bound on the full weighted trace norms \eqref{eq:cstar}.
Those require a valid argument that retains the whole kernel.

\end{document}
